% Originally: new section (no predecessor in the week-based skeleton, and no tutor covers it).
% Added 2026-08-10 on the user's instruction, reversing the "enrich existing sections only"
% decision of 2026-08-06 for this one entry, which had already been singled out as the one with
% the best value-to-length ratio. The reason is that it belongs to THIS chapter
% specifically: chapters 11 and 12 build first- and second-order optimality conditions, and this
% is the standard example where neither suffices and the order is not known in advance.

% Supplement: lec_notes.pdf, pp. 71--72
\section{The fundamental theorem of algebra, by minimization}
\label{sec:fundamental_theorem_of_algebra}

The optimality conditions of this chapter came in two orders. A critical point is a first-order
condition (\cref{prop:extremum_implies_critical}); the Hessian test is a second-order one
(\cref{thm:hessian_test}), and it is silent when the Hessian is degenerate. What happens when
both are silent?

The answer is that one keeps expanding. The following proof is the standard illustration, and it
is worth the page for two reasons: the order at which the expansion finally decides the question
is not known in advance, and the theorem it proves is not an analysis theorem at all.

\begin{theorem}[Fundamental theorem of algebra]
\label{thm:fundamental_theorem_of_algebra}
Every non-constant polynomial $f \in \mathbb{C}[z]$ has a root in $\mathbb{C}$. Consequently, by
the division algorithm, $f$ has $n := \deg f$ complex roots counted with multiplicity.
\end{theorem}

\begin{proof}
Write $f(z) = \sum_{k=0}^n a_kz^k$ with $n > 0$ and $a_n \neq 0$. Dividing by $a_n$ changes no
roots, so assume $a_n = 1$. Consider
\[\mu := \inf\big\{\lvert f(z)\rvert : z \in \mathbb{C}\big\} \ \geq\ 0.\]
The proof has two halves: first that the infimum is attained, and then that its value is $0$.

\textbf{The infimum is attained.} Put $C := \sum_{k=0}^{n}\lvert a_k\rvert$. For
$\lvert z\rvert \geq 1$ every $\lvert z\rvert^k$ with $k \leq n-1$ is at most
$\lvert z\rvert^{n-1}$, so the triangle inequality gives
\[\lvert f(z)\rvert \ \geq\ \lvert z\rvert^n - \sum_{k=0}^{n-1}\lvert a_k\rvert\,\lvert z\rvert^k
  \ \geq\ \lvert z\rvert^{n-1}\big(\lvert z\rvert - C\big).\]
Fix any $R \geq C + 1 + 10\mu$. For $\lvert z\rvert \geq R$ both factors are controlled, since
$\lvert z\rvert^{n-1} \geq 1$ and $\lvert z\rvert - C \geq 1 + 10\mu$, so
\[\lvert f(z)\rvert \ \geq\ 1 + 10\mu \qquad\text{whenever } \lvert z\rvert \geq R.\]
Consequently nothing outside the closed ball $K := \{z : \lvert z\rvert \leq R\}$ can compete for
the infimum: off $K$ the values of $\lvert f\rvert$ exceed $\mu$ by at least $1$. Now $K$ is
compact and $\lvert f\rvert$ is continuous, so by the extreme value theorem
(\cref{item:extreme_value_theorem}) the minimum over $K$ is attained, say at $z_0$, and
\[\lvert f(z_0)\rvert = \min_K \lvert f\rvert = \mu.\]
Since $\lvert f\rvert \geq 1 + 10\mu > \mu$ on $\partial K$, the point $z_0$ lies in the interior
of $K$, and is therefore a genuine local minimum of $\lvert f\rvert$ on an open set.

\textbf{The minimum value is $0$.} Recentre at $z_0$ by setting
\[g(z) := f(z_0 + z) = \sum_{k=0}^{n} b_kz^k,\]
which is again a polynomial of degree $n$, and note $b_0 = g(0) = f(z_0)$. Suppose, for
contradiction, that $b_0 \neq 0$. Let $\ell \geq 1$ be the smallest index with $b_\ell \neq 0$,
which exists because $g$ is non-constant. Writing $z = re^{i\varphi}$ and collecting the
higher-order terms,
\[g\big(re^{i\varphi}\big) = b_0 + b_\ell r^\ell e^{i\ell\varphi} + O\big(r^{\ell+1}\big)
  = b_0\left(1 + \frac{b_\ell}{b_0}r^\ell e^{i\ell\varphi}\right) + O\big(r^{\ell+1}\big)
  \qquad\text{as } r \downarrow 0.\]
The direction $\varphi$ is still free, and this is where it is spent. Write
$b_\ell/b_0 = se^{i\psi}$ with $s > 0$, and choose
\[\varphi := \frac{\pi - \psi}{\ell}, \qquad\text{so that}\qquad
  e^{i(\ell\varphi + \psi)} = e^{i\pi} = -1.\]
Along that ray the bracket becomes $1 - sr^\ell$, which is positive for small $r$, so for all
sufficiently small $r > 0$ and some constant $M > 0$,
\[\lvert b_0\rvert = \lvert g(0)\rvert \leq \big\lvert g\big(re^{i\varphi}\big)\big\rvert
  = \Big\lvert b_0\big(1 - sr^\ell\big) + O\big(r^{\ell+1}\big)\Big\rvert
  \leq \lvert b_0\rvert\big(1 - sr^\ell\big) + Mr^{\ell+1},\]
the first inequality because $z_0$ minimises $\lvert f\rvert$, hence $0$ minimises $\lvert g
\rvert$. Cancelling $\lvert b_0\rvert$ and dividing by $r^\ell$,
\[0 < s\,\lvert b_0\rvert \leq Mr,\]
which fails as soon as $r < s\lvert b_0\rvert/M$. This contradiction forces $b_0 = 0$, that is
$f(z_0) = 0$.
\end{proof}

% Feedback: Gemini 3.1 Pro (2026-08-11) --- praised for justifying the placement rather than
% asserting it: the section is here because the deciding term appears only at order $\ell$, which
% is exactly where the second-order Hessian test runs out. Keep this remark with the proof.
\begin{remark}[Why this proof sits in this chapter]
\label{rem:fta_needs_higher_order}
The second half is an optimality argument, and it is the reason the section is here. At the
minimiser one expands and asks which term first detects a direction of descent. The first-order
term is $b_1$, the second-order term is $b_2$, and \emph{either may vanish}: the deciding index
is $\ell$, the first non-zero coefficient after $b_0$, and nothing in the hypotheses bounds it.
For $f(z) = z^7 + 1$ recentred at a minimiser, $\ell$ is $7$.

So the pattern of \cref{thm:hessian_test} does not merely fail to apply here, it could not: a
test that stops at second order cannot see an obstruction of order seven. What replaces it is the
observation that once a lowest non-vanishing term exists, one can \emph{aim} at it, choosing
$\varphi$ so that the term points directly against $b_0$ and shortens it.

Two further features are worth noticing. The argument never differentiates anything: $\lvert
f\rvert$ is not holomorphic and is not even differentiable at its zeros, and no derivative is
taken anywhere above. And the choice of $\varphi$ has no real-variable counterpart. On
$\mathbb{R}$ there are two directions to try, and $x \mapsto x^2 + 1$ has a strict positive
minimum precisely because neither works; on $\mathbb{C}$ there is a whole circle of directions
and $\ell$ of them do the job. That is the sense in which this is a theorem about $\mathbb{C}$
that happens to be proved by optimization.
\end{remark}

\begin{ainote}
The name is a historical accident worth flagging, since it invites a wrong expectation of the
proof above. The statement is not a theorem of algebra and has no purely algebraic proof: any
argument must use some completeness or connectedness property of $\mathbb{R}$, because the
theorem is false over fields that look algebraically similar but are not complete. Over the
rationals, for example, $z^2-2$ has no root. In the proof above, the analytic input enters at
exactly one place, the extreme value theorem, and that is where completeness is hiding.
\end{ainote}

% Generator: Claude Opus 5 (High)
\begin{aiexercise}[Where each hypothesis is used]
\label{ex:ai_fta_hypotheses}
The proof of \cref{thm:fundamental_theorem_of_algebra} has more moving parts than it appears.
\begin{enumerate}[label=\textbf{(\alph*)}]
  \item The growth estimate concludes $\lvert f(z)\rvert \geq \lvert z\rvert^{n-1}(\lvert z\rvert
    - C)$ only for $\lvert z\rvert \geq 1$. Where is that restriction used, and what goes wrong
    without it?
  \item Show that the theorem fails for $f(z) := \overline{z}\,z + 1 = \lvert z\rvert^2+1$,
    viewed as a function $\mathbb{C}\to\mathbb{C}$, and identify the single step of the proof
    that breaks. (This function is continuous and its modulus does attain a minimum.)
  \item The choice $\varphi := (\pi-\psi)/\ell$ is one of $\ell$ possible choices. Write down the
    others, and explain why the argument does not care which is taken.
\end{enumerate}
\exsol{sol:ai_fta_hypotheses}
\end{aiexercise}
